\subsection{Harmonic synthesis}

Let A be a commutative Banach algebra. Recall that if M is a subset of \(\mathsf{X}(\mathsf{A})\) , we denote by \(\Upsilon(\mathsf{M})\) the intersection of the kernels of the elements of M (cf.~I, p.~30); this is an ideal of A.

PROPOSITION 4. --- Let A be a regular commutative Banach algebra without radical. Let F be a closed subset of X(A). The set of ideals I of A such that V(I) = F, ordered by inclusion, admits a greatest element, namely Y(F), and a least element, namely the set J of x ∈ A such that G(x) has compact support disjoint from F.

By construction, \(\Upsilon(F)\) is an ideal of A such that \(V(\Upsilon(F))\) contains F, and it is the largest ideal of A having this property. Since F is closed, there exists an ideal I of A such that \(V(I) = F\) ; we therefore have \(I \subset \Upsilon(F)\) , whence \(V(\Upsilon(F)) \subset V(I) = F\) so that \(V(\Upsilon(F)) = F\) . This proves the first assertion.

The set J is an ideal of A and V(J) contains F. Let us show that V(J) = F. Let \(\chi \in \mathsf{X}(\mathsf{A})\) not belong to F. Let U be a compact neighborhood of \(\chi\) not meeting F (TG, I, p.~65, cor. of prop. 9). By assertion (ii) of prop. 1 of I, p.~88, there exists \(x \in A\) such that \(\mathcal{G}(x)\) is equal to 1 at \(\chi\) and to 0 outside of U. We then have \(x \in J\) and therefore \(\chi \notin \mathrm{V}(\mathrm{J})\) . This shows that V(J) ⊂ F and hence V(J) = F.

Finally, let I be an ideal of A such that \(V(I) = F\) . We show that \(J \subset I\) . Let \(x \in J\) and let C be the support of \(\mathcal{G}(x)\) ; the set C is a compact subset of \(X(A)\) disjoint from F. By Prop. 3, there exists an element \(u \in I\) such that \(\mathcal{G}(u) = 1\) on C. We then have \(\mathcal{G}(x) = \mathcal{G}(ux)\) and therefore x = ux since A is radical-free (prop. 8 of I, p.~38). Consequently, we have \(x \in I\) , which shows that \(J \subset I\) .

COROLLARY 1. --- Let A be a regular commutative Banach algebra without radical. Let J be the set of \(x \in \mathrm{A}\) such that \(\mathcal{G}(x)\) has compact support. We assume that \(\overline{\mathrm{J}} = \mathrm{A}\) . Then every closed ideal of A distinct from A is contained in a regular maximal ideal.

If I is a closed ideal of A which is contained in no regular maximal ideal, then \(V(I) = \varnothing\) , hence \(I \supset J\) (prop. 4 applied to \(F = \varnothing\) ), whence \(I \supset \overline{J} = A\) .

COROLLARY 2. --- Let A be a commutative regular Banach algebra without radical. Let \(x, y \in A\) . If the support of \(\mathcal{G}(x)\) is compact and contained in the set of characters \(\chi\) such that \(\mathcal{G}(y)(\chi) \neq 0\) , then \(x\) is a multiple of \(y\) in \(A\) .

Let I be the ideal Ay of A. Then V(I) is the set of zeros of \(\mathcal{G}(y)\) . Since the support F of \(\mathcal{G}(x)\) is compact and disjoint from V(I), we have \(x \in I\) (prop. 4 applied to F).

DEFINITION 2. --- Let A be a commutative Banach algebra.

Let I be an ideal of A, \(x \in A\) , and \(\chi \in \mathsf{X}'(\mathsf{A})\) . We say that x belongs to I in a neighborhood of \(\chi\) if there exists an element \(y \in I\) such that \(\mathcal{G}'(y)\) and \(\mathcal{G}'(x)\) coincide in a neighborhood of \(\chi\) .

We say that A satisfies the Ditkin condition if, for every \(\chi \in \mathsf{X}'(\mathrm{A})\) and every \(x \in \mathrm{A}\) such that \(\mathcal{G}'(x)\) vanishes at \(\chi\) there exists a sequence \((x_n)\) in A such that \(x = \lim_{n \to \infty} x_n x\) and such that each \(\mathcal{G}'(x_n)\) vanishes in a neighborhood \(\mathrm{V}_n\) of \(\chi\) .

Remarks. --- Let A be a commutative Banach algebra.

\begin{enumerate}
\def\labelenumi{\arabic{enumi})}
\item
  If \(\chi\) is such that \(\mathcal{G}'(x)\) vanishes in a neighborhood of \(\chi\) , then \(x\) belongs to I in a neighborhood of \(\chi\) .
\item
  If x belongs to I in a neighborhood of \(\chi\) and \(y \in A\) is an arbitrary element, then xy belongs to I in a neighborhood of \(\chi\) .
\item
  The set of \(\chi\) such that \(x\) belongs to I in a neighborhood of \(\chi\) is open in \(\mathrm{X}'(\mathrm{A})\) .
\item
  Suppose that A is regular and radical-free. If x belongs to I in a neighborhood of \(\chi\) for all \(\chi \in \mathsf{X}'(\mathrm{A})\) , then x belongs to I (cor. 2 of I, p.~91 applied to the function \(f = \mathcal{G}'(x)\) and prop. 8 of I, p.~38).
\item
  Suppose that A is regular. Let I be an ideal of A and let \(\chi\) be an element of \(\mathsf{X}(\mathsf{A})\) such that \(\chi\notin\mathsf{V}(\mathsf{I})\) . Then every element \(x\) of A belongs to I in the neighborhood of \(\chi\) . Indeed, by def. 1 of I, p.~89, there exists a \(z\in\mathsf{A}\) such that \(\mathcal{G}'(z)\) is equal to 1 in a neighborhood of \(\chi\) and equal to 0 in a neighborhood of V(I). The support of \(\mathcal{G}(z)\) is compact and hence we have \(z\in\mathsf{I}\) (prop. 4 applied to V(I)); therefore \(xz\in\mathsf{I}\) , and \(\mathcal{G}'(xz)=\mathcal{G}'(x)\) in a neighborhood of \(\chi\) .
\end{enumerate}

Recall that a subspace K of a topological space X is said to be perfect if it is closed and has no isolated point (TG, I, p.~8).

Lemma 1. --- Let A be a commutative regular Banach algebra without radical, satisfying Ditkin's condition. Let I be a closed ideal of A and let x be an element of \(\Upsilon(\mathrm{V}(\mathrm{I}))\) . Let K be the set of \(\chi \in \mathsf{X}'(\mathrm{A})\) such that x does not belong to I in a neighborhood of \(\chi\) . Then the set K is a perfect subset of \(\mathsf{X}'(\mathrm{A})\) .

Let G denote the complement of K in \(X'(A)\) . The set G is open in \(X'(A)\) (remark 3), hence K is closed.

Proceed by contradiction, and suppose that K has an isolated point \(\chi_{0}\) . Let U be a neighborhood of \(\chi_{0}\) such that \(U - \{\chi_{0}\}\subset G\) . Since x does not belong to I in a neighborhood of \(\chi_{0}\) , Remark 5 shows that \(\chi_{0}\in\mathrm{V(I)}\) . In particular, we have \(\chi_{0}(x)=0\) since \(x\in\Upsilon(\mathrm{V(I)})\) .

We are going to show that there exists an element \(y\) of A which belongs to I in a neighborhood of every point of \(\mathsf{X}'(\mathrm{A}) - \{\chi_0\}\) , which does not belong to I in a neighborhood of \(\chi_0\) , and such that \(\chi_0(y) = 0\) .

Suppose first that the existence of such an element \(y\) has been proved. Since A satisfies Ditkin's condition, there then exists a sequence \((x_n)\) in A such that \(x_n y\) tends to \(y\) and such that each \(\mathcal{G}'(x_n)\) vanishes in a neighborhood of \(\chi_0\) . For every \(n\) , the element \(x_n y\) then belongs to I in a neighborhood of every point of \(X'(A)\) (Remarks 1 and 2) and therefore \(x_{n}y \in I\) (remark 4). Since I is closed, we deduce that \(y \in I\) , which contradicts the fact that y does not belong to I in a neighborhood of \(\chi_{0}\) .

It remains to prove the existence of \(y\) . If \(\chi_0 \neq 0\) , by assertion (iii) of prop. 1 of I, p.~88, there exists a \(u \in A\) such that \(\mathcal{G}'(u)\) is equal to 1 in a neighborhood of \(\chi_0\) and equal to 0 in a neighborhood of \(\mathrm{X}'(\mathrm{A}) - \mathrm{U}\) . Let \(y = ux\) . Since \(x\) belongs to I in a neighborhood of \(\chi\) for all \(\chi \in \mathrm{U} - \{\chi_0\}\) , the same is true of \(y\) . Moreover, if \(\chi \in \mathrm{X}'(\mathrm{A}) - \mathrm{U}\) , then \(\mathcal{G}'(y)\) vanishes in a neighborhood of \(\chi\) . Therefore (Remark 5) the element \(y = ux\) belongs to I in a neighborhood of every \(\chi \neq \chi_0\) . Since \(\mathcal{G}'(y)\) coincides with \(\mathcal{G}'(x)\) in a neighborhood of \(\chi_0\) , the fact that \(\chi_0\) belongs to K implies that \(y\) does not belong to I in a neighborhood of \(\chi_0\) . Finally, we have \(\chi_0(y) = \chi_0(u)\chi_0(x) = 0\) .

If \(\chi_0 = 0\) , there similarly exists an element \(v \in \mathrm{A}\) such that \(\mathcal{G}'(v)\) is zero in a neighborhood of \(\chi_0\) and equal to 1 in a neighborhood of \(\mathsf{X}'(\mathsf{A}) - \mathsf{U}\) ; as before, we deduce from this that the element \(y = x - vx\) belongs to I in a neighborhood of every \(\chi \neq \chi_0\) , that it does not belong to I in a neighborhood of \(\chi_0\) , and that \(\chi_0(y) = 0\) .

Lemma 2. — Let X be a topological space. Let F and D be disjoint subspaces of X such that F is closed and D is discrete. If F contains no nonempty perfect subspace, then neither does F ∪ D.

Indeed, suppose that K is a nonempty perfect subspace of F ∪ D. Let x be a point of K. If x belongs to D, it is isolated in D, hence also in F ∪ D since F is closed. Therefore x is isolated in K, which contradicts the hypotheses.

PROPOSITION 5. — Let A be a regular commutative Banach algebra without radical, satisfying Ditkin's condition. Let I be a closed ideal of A such that the boundary F of V(I) contains no nonempty perfect set. Then I = Y(V(I)), that is, I is the set of x ∈ A such that G(x) vanishes on V(I). In particular, if V(I) reduces to a single point χ, we have I = Ker(χ).

We have \(\mathrm{I} \subset \Upsilon(\mathrm{V}(\mathrm{I}))\) . Now let \(x \in \Upsilon(\mathrm{V}(\mathrm{I}))\) . Let \(\mathrm{G}\) be the set of characters \(\chi \in \mathsf{X}'(\mathbf{A})\) such that \(x\) belongs to \(\mathrm{I}\) in a neighborhood of \(\chi\) . It is open and its complement \(\mathrm{K}\) is perfect (lemma 1). Since \(\mathcal{G}'(x)\) is zero on \(\mathrm{V}(\mathrm{I})\) , the set \(\mathrm{G}\) contains the interior of \(\mathrm{V}(\mathrm{I}) \cup \{0\}\) (remark 1). It also contains \(\mathsf{X}(\mathbf{A}) - \mathsf{V}(\mathbf{I})\) by remark 5. Therefore \(\mathrm{K} = \mathsf{X}'(\mathbf{A}) - \mathsf{G}\) is contained in the boundary \(\mathrm{F}_0\) of \(\mathrm{V}(\mathrm{I}) \cup \{0\}\) .

We have \(F_{0} \subset F \cup \{0\}\) . The hypothesis therefore implies that \(F_{0}\) contains no nonempty perfect set (lemma 2). It follows therefore from lemma 1 that the perfect set K is empty. Hence x belongs to I in a neighborhood of every \(\chi \in \mathsf{X}'(\mathsf{A})\) , which means that \(x \in I\) (remark 4). Thus \(\Upsilon(\mathrm{V}(\mathrm{I})) \subset \mathrm{I}\) , which concludes the proof.

